\documentclass[10pt]{article}
%\usepackage[T1]{fontenc}
\usepackage[utf8]{inputenc}
%\usepackage[francais]{babel}
%\usepackage{listingsutf8}

\usepackage[french]{babel}
%\usepackage[T1]{fontenc}
\usepackage{comment,fullpage}
\usepackage{color}
\usepackage{graphicx}
\usepackage{lastpage}
\usepackage{amsmath}

\setlength{\parindent}{0pt}
\setlength{\parskip}{3pt}

\newcommand{\red}[1]{\textcolor{red}{#1}}

%\newenvironment{solution}{\paragraph{solution}}{}
\newenvironment{solution}{\vspace{6pt}\noindent\color{red}\textbf{solution : }}{\vspace{6pt}\ignorespacesafterend}

\includecomment{comment}
%\excludecomment{solution}

% MATH -----------------------------------------------------------
\newcommand{\norm}[1]{\left\Vert#1\right\Vert}
\newcommand{\abs}[1]{\left\vert#1\right\vert}
\newcommand{\set}[1]{\left\{#1\right\}}
\newcommand{\Real}{\mathbb R}
\newcommand{\eps}{\varepsilon}
\newcommand{\To}{\longrightarrow}
\newcommand{\BX}{\mathbf{B}(X)}
\newcommand{\A}{\mathcal{A}}
% ----------------------------------------------------------------

\begin{document}

\begin{center}
{\Large NFP 107 -- Systèmes de gestion de bases de données}\\\vspace{1ex}
{\Large Examen septembre 2018. Durée: 2h30.}\\

\end{center}

\vspace{1ex}
\hrule
\vspace{1ex}
\textbf{Consignes :}\\
-- Tous les documents sont autorisés.\\
-- Les téléphones mobiles et autres équipements communicants (exemple: PC, tablette, etc.)  doivent être éteints et rangés dans les sacs pendant toute la durée de l'épreuve.\\
-- Pensez à reporter votre numéro d'anonymat sur toutes les feuilles.\\
-- Merci de soigner votre écriture et de ne pas rédiger au crayon papier sur votre copie.\\
Ce sujet comporte 3 pages d'énoncé.

\vspace{1ex}
\hrule
\vspace{1ex}

%\input{./titlepage.tex}
% \addtocounter{page}{1}

\vspace{1ex}

Le schéma de base de données suivant sera utilisé pour l'ensemble du
sujet. Il permet de gérer la gestion des emprunts de livres pour une bibliothèque :

\begin{verbatim}
LIVRE(id, titre, idAuteur, idEditeur, prix,
      dateAchat, datePublication, quantitéEnStock)
AUTEUR(id, nom, prenom, nationalité)
EDITEUR(id, nom, adresse)
CLIENT(id, nom, prenom, adresse, dateNaissance, nationalité, tel)
EMPRUNT(idClient, idLivre, dateDebut, dateFin, dateRetour)
\end{verbatim}

La date prévue de retour est donné par l'attribut \texttt{EMPRUNT.dateFin}. L'attribut \texttt{EMPRUNT.retour} représente la date réelle de retour et vaut NULL si le livre n'a pas été encore retourné.

\section{Algèbre relationnelle et SQL (9 points)}

\begin{enumerate}

\item (Algèbre et SQL) Noms et prénoms des auteurs qui ont publié un livre après 01/01/2018. (1 point)

\item (Algèbre et SQL) Noms et prénoms des auteurs qui ont publié un livre à la même date. (2 points)

\item (Algèbre et SQL) Noms et prénoms des clients (sans doublons, en ordre alphabétique) qui ont emprunté au moins une fois un livre d'Amelie Nothomb publié chez Albin Michel. (2 points)

\item (Algèbre et SQL) Noms des clients qui n'ont jamais retourné un livre à temps. (1 point)

\item (SQL seulement) Les nationalités des clients triés en ordre décroissant par le nombre de dépassements de la date de retour. (2 points)

\item (Algèbre seulement) Clients qui ont emprunté tous les livres d'Albert Camus. (1 point)

\begin{solution}

\end{solution}

\end{enumerate}

\section{Indexation (4 points)}

\begin{enumerate}

\item Quelles sont, parmi les propriétés suivantes, celles qui caractérisent un arbre B+? (1,5 points)

\begin{enumerate}
  \item La hauteur est fixe.
  \item Tous les chemins de la racine vers une feuille ont la même longueur.
  \item L'arbre se réorganise dynamiquement en fonction des insertions et suppressions dans la table.
  \item Une clé est un ensemble de valeurs $[v_1, v_2, \cdots v_n]$ et l'index 
      permet de faire des recherches sur n'importe quel $v_i$.
  \item Il n'y a qu'un seul arbre B+ par table, sur la clé primaire.
  \item Il y a autant d'entrées dans les feuilles que d'enregistrements dans la table indexée.
\end{enumerate}

\begin{solution}
Réponses 2, 3, 6.
\end{solution}

\item  Quelle est la condition pour créer un index non dense sur un fichier contenant 
un ensemble d'enregistrements (1,5 points)?

\end{enumerate}
\section{Optimisation (3 points)}

\begin{enumerate}
    \item Donnez le plan d'exécution pour une requête au choix parmi celles de la Section~1 (1,5 pts)
    \item Expiquez en quoi la taille de la mémoire peut influer sur l'algorithme de jointure.
    

\begin{solution}
Si une des tables tient entièrement en mémoire, il suffit de la charger dans une table de hachage 
et de faire défiler l'autre table. On obtient un algorithme simple qui lit exactement une fois chaque table.
\end{solution}
    
\end{enumerate}

\section{Concurrence (4 points)}


Soit l'exécution concurrente suivante

$$H = r_1[x] r_2[x] r_3[y] w_1[x] w_1[z] c_1 w_2[y] c_2 w_3[z] c_3$$

\begin{enumerate}

\item Donner la liste des conflits de H; \texttt{(0,5 point)}
\begin{solution}
    \begin{itemize}
        \item[sur x : ] $r_2[x] w_1[x]$
        \item[sur y : ] $r_3[y] w_2[y]$
        \item[sur z : ] $w_1[z] w_3[z]$
    \end{itemize}
\end{solution}

\item Donner le graphe de s\'erialisation de H. Que pouvez-vous déduire de ce
graphe? \texttt{(0,5 point)}
\begin{solution}
    $T_2 \rightarrow T_1, T_3 \rightarrow T_2, T_1 \rightarrow T_3$. Il y a un cycle, dont H n'est pas \textit{s\'erialisable}.
\end{solution}


\item Donner l'histoire finale par application de l'algorithme de verrouillage \`a deux phases. 
Donner le d\'etail du d\'eroulement de l'algorithme.
 \texttt{(2 points)}

\begin{solution}
$T_1$  est mis en attente au moment où $w_1[x]$ est soumis ($x$ a un  verrou partagé avec $T_2$)
$T_2$  est mis en attente au moment où $w_2[y]$ est soumis ($x$ a un  verrou partagé avec $T_3$). 
$T_3$ finit alors de s'exécuter, $T_2$ est libérée et termine, puis $T_1$.
\end{solution}


\end{enumerate}
\end{document}

